\documentclass[11pt]{article}

\usepackage[margin=1in]{geometry}
\usepackage{amsmath, amssymb, amsthm}
\usepackage{bm}
\usepackage{graphicx}
\usepackage{hyperref}
\usepackage{enumitem}
\usepackage{mathtools}
\usepackage{bbm}
\usepackage{listings}
\usepackage{courier}

\lstdefinestyle{code}{
  basicstyle=\footnotesize\ttfamily,
  breaklines=true,
  columns=fullflexible,
  frame=single
}

\title{%
  CROF-LC on Multiplicity: \\
  A Contractive, Prime-Indexed Operator Substrate \\
  for Clinical and Sovereignty-Compliant Dynamics
}
\author{Citizen Gardens \\ \\ The Foundation of Multiplicity}
\date{April 5, 2026}

\begin{document}

\maketitle

\begin{abstract}
This report synthesizes the development of a contractive, prime-indexed operator framework (CROF-LC) embedded in a broader Multiplicity Theory stack.
We formalize the Contractive Recursive Operator Framework with Lawful Channels (CROF-LC), specify a practical learning algorithm for its components, and articulate how its mathematical guardrails subsume conventional regulatory constraints (FDA, HIPAA, SOC~II, ISO) as architectural invariants rather than bolt-on policies.
The document is intended as a bridge between Multiplicity-theoretic ontology, implementable machine learning modules, and clinical/regulatory deployment.
\end{abstract}
\newpage
\tableofcontents
\newpage
\section{Executive Summary}

\subsection{Context and Motivation}

The core idea is to treat clinical and governance AI not as a collection of loosely constrained models, but as a \emph{lawful substrate} whose dynamics are constrained by contraction, prime-indexed multiplicity, and explicit auditability.
Rather than adding safety, ethics, or regulatory compliance as exogenous policies, these properties are encoded into the very form of the evolution operator.

Concretely, we work with a state space
\[
  \mathcal{H} = \bigoplus_{m=1}^{M} \mathcal{H}_m,
\]
a direct sum of Hilbert spaces for different modalities (e.g.\ genomics, biosignals, labs, text).
Over this state space, we define a prime-indexed operator field
\[
  \mathcal{F}(x) = \sum_{p \in P} w_p(x)\,\Pi_p\,\mathcal{O}_p(x),
\]
where each prime channel \(p\) carries an operator \(\mathcal{O}_p\), a stability projector \(\Pi_p\), and a resonance-gated weight \(w_p(x)\).
The global evolution is then
\[
  x_{t+1} = x_t + \Lambda_m(t)\,T(\mathcal{F}(x_t)),
\]
with explicit Lipschitz/contraction constraints.

\subsection{Key Innovations}

\begin{itemize}[leftmargin=*]
  \item \textbf{Prime-Indexed Operator Field:}
  Primes index irreducible operator channels \(\mathcal{O}_p\) acting on a shared multiplicity Hilbert space, extending prime indexing from encoding to dynamics.
  \item \textbf{Contraction as a First-Class Constraint:}
  The global map \(\mathcal{F}\) is enforced to be contractive,
  \(\|\mathcal{F}(x) - \mathcal{F}(y)\| \leq \kappa \|x-y\|\) with \(\kappa < 1\),
  via architectural choices and penalties.
  This yields stability, drift bounds, and unique fixed points under mild conditions.
  \item \textbf{Resonance-Modulated Weights:}
  Weights \(w_p(x)\) depend on cross-modal resonance between, for example, genomic and biosignal embeddings, allowing channels to activate when modalities cohere.
  \item \textbf{Prime Channel Specialization:}
  Sparsity and diversity regularization ensure that only a few prime channels are active per sample and that channels are non-redundant, producing interpretable, disentangled structure.
  \item \textbf{Compliance-by-Construction:}
  Contraction, logging, and traceable state evolution make standard regulatory concerns (FDA CDS guidance, HIPAA flows, SOC~II/ISO controls) natural consequences of the architecture rather than separate layers.
\end{itemize}

\subsection{High-Level Learning Problem}

Instead of pure loss minimization, we consider a constrained identification problem:
\[
  \min_{\{\mathcal{O}_p, \Pi_p, M_p\}} \mathcal{E}(x)
  \quad \text{s.t.} \quad
  \rho(\mathcal{F}) < 1,
\]
where \(\rho\) denotes a Lipschitz/spectral bound, and \(\mathcal{E}\) is an energy functional combining task performance, contraction, calibration, and lawfulness terms.

\section{Formal Definition of CROF-LC}

\subsection{State Space and Modality Decomposition}

Let \(\mathcal{H}\) be a real or complex Hilbert space decomposed as
\[
  \mathcal{H} = \bigoplus_{m=1}^{M} \mathcal{H}_m,
\]
with each \(\mathcal{H}_m\) corresponding to a modality (genomics, biosensors, clinical text, etc.).
We write elements \(x \in \mathcal{H}\) as concatenations \(x = (x^{(1)}, \dots, x^{(M)})\).

\subsection{Prime-Indexed Operator Channels}

Let \(P\) be a finite set of primes, \(|P| = K\).
For each \(p \in P\), define an operator
\[
  \mathcal{O}_p : \mathcal{H} \to \mathcal{H}
\]
with the factorization
\[
  \mathcal{O}_p = A_p \circ B_p \circ C,
\]
where:
\begin{itemize}[leftmargin=*]
  \item \(C : \mathcal{H} \to \mathcal{H}\) is a shared linear encoder (e.g.\ spectrally normalized linear map).
  \item \(B_p : \mathcal{H} \to \mathcal{H}\) is a prime-specific low-rank mixer.
  \item \(A_p : \mathcal{H} \to \mathcal{H}\) is a 1-Lipschitz non-linear block.
\end{itemize}

A concrete linear implementation of \(B_p\) is
\[
  B_p(x) = U_p \Sigma_p V_p^\top x,
\]
where \(U_p, V_p \in \mathbb{R}^{d \times r}\) and \(\Sigma_p \in \mathbb{R}^{r \times r}\) is diagonal, with rank \(r \ll d\).
The spectral norm of \(B_p\) is controlled via the diagonal entries of \(\Sigma_p\).

\subsection{Stability Projectors}

Each prime channel is equipped with a stability projector
\[
  \Pi_p : \mathcal{H} \to \mathcal{H},
\]
implemented in minimal form as a diagonal mask
\[
  \Pi_p = \mathrm{diag}(s_p), \quad s_p \in [0,1]^d,
\]
with an \(\ell_1\) penalty to encourage sparsity in \(s_p\).
This is a simple but effective surrogate for a learned spectral projector onto stable modes.

\subsection{Resonance and Gating}

Let \(x^{(\mathrm{gen})}\) and \(x^{(\mathrm{bio})}\) denote two modality-specific embeddings (e.g.\ genomic and biosensor components).
For each prime \(p\), define a resonance matrix \(M_p\) and score
\[
  R_p(x) = \frac{\langle x^{(\mathrm{gen})}, M_p x^{(\mathrm{bio})} \rangle}{\|x^{(\mathrm{gen})}\|\,\|M_p x^{(\mathrm{bio})}\| + \varepsilon},
\]
with \(\varepsilon > 0\) to avoid division by zero.
The corresponding gate is
\[
  w_p(x) = \sigma\big(R_p(x) - \tau_p\big),
\]
where \(\sigma\) is the logistic sigmoid and \(\tau_p\) is a learned threshold.
This yields a fail-closed pattern: channels with sub-threshold resonance remain effectively inactive.

\subsection{Global Operator and Update Rule}

The prime-indexed operator field is
\[
  \mathcal{F}(x) = \sum_{p \in P} w_p(x)\,\Pi_p\,\mathcal{O}_p(x).
\]
Given a modality-dependent step-size \(\Lambda_m(t)\) and a transfer function
\(T : \mathcal{H} \to \mathcal{H}\) (e.g.\ mapping to risk scores or actions), we define the recursion
\[
  x_{t+1} = x_t + \Lambda_m(t)\,T(\mathcal{F}(x_t)).
\]

\subsection{Contraction Condition}

We require that \(\mathcal{F}\) be globally Lipschitz with constant \(\kappa < 1\):
\[
  \|\mathcal{F}(x) - \mathcal{F}(y)\|
  \leq \kappa \|x-y\|
  \quad \forall x,y \in \mathcal{H}.
\]
Architecturally, this is enforced via:
\begin{itemize}[leftmargin=*]
  \item Spectral normalization of linear layers in \(C\) and the linear constituent of \(A_p\).
  \item Norm constraints on \(B_p\) through singular values \(\Sigma_p\).
  \item Control of the effective weight budget \(\sum_p \|w_p\|_{\infty}\).
\end{itemize}
A design-time bound of the form
\[
  \|\mathcal{F}\|_{\mathrm{Lip}} \leq W_{\max} \cdot \alpha < 1
\]
can be guaranteed by choosing \(\alpha\) and \(W_{\max}\) appropriately and monitoring them during training.

\section{Learning Problem and Energy Functional}

\subsection{Task and Energy}

Let \(f_\Theta\) denote the CROF-LC-based predictor with parameters
\(\Theta := \{\mathcal{O}_p, \Pi_p, M_p\}_{p \in P}\).
Given data \((x,y)\) and a predictive output \(\hat{y} = f_\Theta(x)\), we consider a total loss
\begin{align*}
  \mathcal{L}_{\text{total}} &=
  \underbrace{\mathcal{L}_{\text{task}}(\hat{y}, y)}_{\text{prediction}}
  + \lambda_1 \underbrace{\|\mathcal{F}(x) - x\|^2}_{\text{contraction energy}} \\
  &\quad + \lambda_2 \underbrace{H(p(\hat{y}\mid x))}_{\text{entropy penalty}}
  + \lambda_3 \underbrace{D_{\mathrm{KL}}(p(\hat{y}\mid x)\,\|\,\text{prior})}_{\text{calibration}} \\
  &\quad + \lambda_4 \underbrace{\mathcal{L}_{\mathrm{Lip}}}_{\text{stability constraint}}
  + \lambda_5 \underbrace{\mathcal{L}_{\mathrm{sparsity}}}_{\text{prime sparsity}} \\
  &\quad + \lambda_6 \underbrace{\mathcal{L}_{\mathrm{div}}}_{\text{prime diversity}}
  + \lambda_7 \underbrace{\mathcal{L}_{\mathrm{proj}}}_{\text{projector sparsity}}
  + \lambda_8 \underbrace{\mathcal{L}_M}_{\text{resonance regularization}}.
\end{align*}

Typical choices:
\begin{itemize}[leftmargin=*]
  \item \(\mathcal{L}_{\text{task}}\): cross-entropy, Brier score, survival loss, etc.
  \item \(H\): Shannon entropy of the predictive distribution.
  \item \(D_{\mathrm{KL}}\): KL divergence to a clinically calibrated prior distribution.
\end{itemize}

\subsection{Sparsity and Diversity Terms}

Let \(z_p(x) = \Pi_p \mathcal{O}_p(x)\) denote the prime-specific channel output.
Then:
\begin{align*}
  \mathcal{L}_{\mathrm{sparsity}} &= \sum_{p \in P} \mathbb{E}_x \big[ |w_p(x)| \big], \\
  \mathcal{L}_{\mathrm{div}} &= \sum_{\substack{p,q \in P \\ p < q}}
    \mathbb{E}_x \big[ \langle z_p(x), z_q(x) \rangle^2 \big], \\
  \mathcal{L}_{\mathrm{proj}} &= \sum_{p \in P} \|s_p\|_1, \\
  \mathcal{L}_M &= \sum_{p \in P} \|M_p\|_F^2.
\end{align*}

These encourage:
\begin{itemize}[leftmargin=*]
  \item Few active prime channels per sample.
  \item Low redundancy among channels.
  \item Interpretable projector masks.
  \item Regularized resonance matrices.
\end{itemize}

\subsection{Lipschitz Penalty}

In addition to architectural constraints, one may include an explicit Lipschitz penalty:
\begin{itemize}[leftmargin=*]
  \item Empirical global estimate:
  \[
    \hat{\kappa} = \max_{i,j} \frac{\|\mathcal{F}(x_i) - \mathcal{F}(x_j)\|}{\|x_i - x_j\|},
  \]
  with \(\mathcal{L}_{\mathrm{Lip}} = \max(0, \hat{\kappa} - \kappa_{\text{target}})\).
  \item Jacobian spectral norm penalty:
  \[
    \mathcal{L}_{\mathrm{Lip}} = \mathbb{E}_x\big[ \|J_{\mathcal{F}}(x)\|_2 \big].
  \]
\end{itemize}
In practice, a minimal viable implementation can rely solely on spectral normalization and weight-budget control, using Jacobian-based estimates for diagnostics.

\section{Learning Algorithm and Training Phases}

\subsection{Parameterization Summary}

We summarize the minimal parameterization for a practical implementation:

\begin{itemize}[leftmargin=*]
  \item Shared encoder \(C\): linear map with spectral normalization,
  \(\|C\|_2 \leq 1\).
  \item Prime-specific mixer \(B_p\): low-rank map \(U_p \Sigma_p V_p^\top\) with constrained singular values.
  \item Prime-specific non-linearity \(A_p\): 1-Lipschitz block (e.g.\ GroupSort).
  \item Projectors \(\Pi_p\): diagonal masks with entries \(s_p \in [0,1]\).
  \item Resonance matrices \(M_p\): small matrices with Frobenius penalties.
\end{itemize}

\subsection{Phase A: Warm Start (Structure First)}

\begin{itemize}[leftmargin=*]
  \item Spectral normalization active but with loose bounds (\(\kappa\) not yet strongly enforced).
  \item Loss dominated by \(\mathcal{L}_{\text{task}}\) and a small contraction term.
  \item Sparsity, diversity, and explicit Lipschitz penalties either off or very small.
\end{itemize}
Goal: learn a useful representation and non-trivial prime channels without yet demanding strict contraction.

\subsection{Phase B: Contraction Enforcement}

\begin{itemize}[leftmargin=*]
  \item Tighten spectral norm caps and set an initial Lipschitz target \(\kappa_{\text{target}} \approx 1\).
  \item Gradually anneal \(\kappa_{\text{target}}\) to the desired value (e.g.\ \(0.9 \to 0.7\)).
  \item Introduce or increase \(\lambda_4\) for \(\mathcal{L}_{\mathrm{Lip}}\).
\end{itemize}

\subsection{Phase C: Prime Channel Specialization}

\begin{itemize}[leftmargin=*]
  \item Increase \(\lambda_5\) and \(\lambda_6\) to promote sparsity and diversity.
  \item Monitor metrics such as the average number of active primes per sample and pairwise correlations of \(z_p\).
\end{itemize}

\subsection{Phase D: Calibration and Lawfulness}

\begin{itemize}[leftmargin=*]
  \item Increase \(\lambda_2, \lambda_3\) to improve entropy and KL-based calibration.
  \item Optionally introduce additional lawfulness energy terms (e.g.\ enforcing projections onto lawful subspaces).
  \item Freeze most of the operator structure and fine-tune thresholds \(\tau_p\) and small heads for calibration.
\end{itemize}

\section{Information-Theoretic Extension}

To formalize the emergence of prime channels, we may add an explicit mutual-information / total-correlation objective.

\subsection{Mutual Information with Targets}

Let \(A\) denote the random variable representing the active prime pattern (for example via binary or continuous activations).
We can consider an auxiliary term
\[
  \mathcal{L}_{\mathrm{MI}} = -I(A; Y),
\]
encouraging prime activation patterns to be informative about labels \(Y\).

\subsection{Redundancy via Total Correlation}

Let \(Z_1, \dots, Z_K\) denote the prime channel outputs \(z_p(x)\).
The total correlation (multi-information) is
\[
  \mathrm{TC}(Z_1,\dots,Z_K) =
    \sum_{p=1}^{K} H(Z_p) - H(Z_1,\dots,Z_K).
\]
A regularizer
\[
  \mathcal{L}_{\mathrm{TC}} = \mathrm{TC}(Z_1,\dots,Z_K)
\]
penalizes redundancy across channels.
An information-theoretic layer can then be
\[
  \mathcal{L}_{\mathrm{info}} = \mathcal{L}_{\mathrm{MI}} + \beta\,\mathcal{L}_{\mathrm{TC}},
\]
with \(\beta > 0\).
In practice, both \(I\) and \(\mathrm{TC}\) can be approximated with variational mutual information estimators.

\section{Alignment with Multiplicity Theory}

\subsection{Prime-Indexed Multiplicity}

In the Multiplicity framework, primes act as labels for irreducible multiplicity classes or channels of lawful recursion.
CROF-LC realizes this by assigning each prime \(p\) an operator family \(\mathcal{O}_p\), projector \(\Pi_p\), and resonance pattern \(w_p(x)\).
The decomposition
\[
  \mathcal{F}(x) = \sum_{p} w_p(x)\,\Pi_p\,\mathcal{O}_p(x)
\]
is thus a concrete multiplicity operator, with the learning process discovering which prime channels become active for which regimes.

\subsection{Contraction and Fixed Points}

The contraction condition and the energy term \(\|\mathcal{F}(x) - x\|^2\) place CROF-LC in the same conceptual space as multiplicity-inspired fixed-point and ACE-style algorithms:
\begin{itemize}[leftmargin=*]
  \item Unique stable fixed points under contraction.
  \item Bounded evolution and controlled drift.
  \item A direct bridge to projected dynamical systems and viability kernels.
\end{itemize}

\subsection{Lawfulness Manifolds and Ethics as Constraints}

By design, CROF-LC restricts the dynamics to a subset of maps satisfying contraction, calibration, and sparsity/diversity constraints.
One can interpret this subset as a ``lawfulness manifold'' in the space of possible evolutions, with ethics and sovereignty encoded as constraints on evolution rather than external preferences.

\section{Illustrative Code Snippet (PyTorch-Style)}

\subsection{CROF-LC Module Skeleton}

\begin{lstlisting}[style=code,language=Python,caption={Skeleton CROF-LC module in PyTorch-style pseudocode}]
import torch
import torch.nn as nn
import torch.nn.functional as F
from torch.nn.utils import spectral_norm

class GroupSort(nn.Module):
    def __init__(self, group_size=2):
        super().__init__()
        self.group_size = group_size

    def forward(self, x):
        b, d = x.shape
        g = self.group_size
        assert d % g == 0
        x = x.view(b, d // g, g)
        x, _ = torch.sort(x, dim=-1)
        return x.reshape(b, d)

class LowRankMixer(nn.Module):
    def __init__(self, dim, rank):
        super().__init__()
        self.U = nn.Parameter(torch.randn(dim, rank) * 0.02)
        self.log_s = nn.Parameter(torch.zeros(rank))
        self.V = nn.Parameter(torch.randn(dim, rank) * 0.02)

    def forward(self, x):
        s = torch.exp(self.log_s).clamp(max=1.0)
        return x @ self.V @ torch.diag(s) @ self.U.T

class PrimeOperator(nn.Module):
    def __init__(self, dim, rank):
        super().__init__()
        self.mixer = LowRankMixer(dim, rank)
        self.lin = spectral_norm(nn.Linear(dim, dim))
        self.act = GroupSort(group_size=2)

    def forward(self, x):
        z = self.mixer(x)
        z = self.lin(z)
        z = self.act(z)
        return z

class ResonanceGate(nn.Module):
    def __init__(self, dim_gen, dim_bio):
        super().__init__()
        self.M = nn.Parameter(torch.randn(dim_bio, dim_gen) * 0.02)
        self.tau = nn.Parameter(torch.zeros(1))

    def forward(self, x_gen, x_bio, eps=1e-6):
        proj = x_bio @ self.M
        num = (x_gen * proj).sum(dim=-1, keepdim=True)
        den = x_gen.norm(dim=-1, keepdim=True) * proj.norm(dim=-1, keepdim=True) + eps
        r = num / den
        return torch.sigmoid(r - self.tau)

class CROFLC(nn.Module):
    def __init__(self, input_dim, hidden_dim, primes, rank):
        super().__init__()
        self.primes = primes
        self.encoder = spectral_norm(nn.Linear(input_dim, hidden_dim))
        self.ops = nn.ModuleDict({
            str(p): PrimeOperator(hidden_dim, rank) for p in primes
        })
        self.projectors = nn.ParameterDict({
            str(p): nn.Parameter(torch.zeros(hidden_dim)) for p in primes
        })
        self.gates = nn.ModuleDict({
            str(p): ResonanceGate(hidden_dim, hidden_dim) for p in primes
        })
        self.head = nn.Linear(hidden_dim, hidden_dim)

    def forward(self, x):
        h = self.encoder(x)
        x_gen, x_bio = h, h
        zs, ws = {}, {}
        agg = 0.0
        for p in self.primes:
            key = str(p)
            z = self.ops[key](h)
            s = torch.sigmoid(self.projectors[key])
            z = z * s
            w = self.gates[key](x_gen, x_bio)
            agg = agg + w * z
            zs[key] = z
            ws[key] = w
        out = self.head(agg)
        return out, zs, ws
\end{lstlisting}

\subsection{Loss Function Sketch}

\begin{lstlisting}[style=code,language=Python,caption={Composite loss with contraction, calibration, sparsity, and diversity}]
def crof_loss(pred, target, x, F_x, zs, ws, prior=None,
              l_con=1.0, l_ent=0.01, l_kl=0.01,
              l_sparse=0.01, l_div=0.01, l_proj=0.01):
    loss_task = F.cross_entropy(pred, target)
    loss_con = ((F_x - x) ** 2).mean()
    probs = pred.softmax(dim=-1)
    loss_ent = -(probs * (probs.clamp_min(1e-8).log())).sum(dim=-1).mean()
    if prior is None:
        prior = torch.full_like(probs, 1.0 / probs.size(-1))
    loss_kl = (probs * (probs.clamp_min(1e-8).log() - prior.clamp_min(1e-8).log())).sum(dim=-1).mean()
    loss_sparse = sum(w.abs().mean() for w in ws.values())
    z_list = list(zs.values())
    loss_div = 0.0
    for i in range(len(z_list)):
        for j in range(i + 1, len(z_list)):
            loss_div = loss_div + ((z_list[i] * z_list[j]).sum(dim=-1) ** 2).mean()
    loss_proj = 0.0
    for z in zs.values():
        loss_proj = loss_proj + z.abs().mean()
    total = (loss_task + l_con*loss_con + l_ent*loss_ent + l_kl*loss_kl
             + l_sparse*loss_sparse + l_div*loss_div + l_proj*loss_proj)
    return total
\end{lstlisting}

\section{Recommended Next Steps}

\begin{enumerate}[leftmargin=*]
  \item Implement CROF-LC with a small number of prime channels (e.g.\ 8–16) and low-rank mixers on a single-modality clinical time series dataset (vitals/labs).
  \item Benchmark against standard baselines (LSTM, Transformer) on accuracy, temporal stability of risk trajectories, and calibration.
  \item Add multi-modal resonance and information-theoretic regularizers to study emergent prime-channel semantics.
  \item Integrate logging of active primes, contraction estimates, and energy components as a proto-audit layer for clinical and regulatory stakeholders.
\end{enumerate}

\section*{Mathematical Appendix}

\addcontentsline{toc}{section}{Mathematical Appendix}

In this appendix we collect explicit operator norm bounds and proof sketches that justify the contraction properties and fixed-point behavior of the CROF-LC operator field, including the effect of resonance maps and gates, and the stability of the full recursion
\[
  x_{t+1} = x_t + \Lambda T(\mathcal{F}(x_t)).
\]

\subsection*{A.1 Lipschitz Bounds for Channel Operators}

We first bound the Lipschitz constants of the prime-channel operators \(\mathcal{O}_p\) and then of the global operator \(\mathcal{F}\), in a way that is compatible with our architectural choices.

\begin{theorem}[Channel Lipschitz Bound]
Assume:
\begin{enumerate}[label=(\roman*),leftmargin=*]
  \item \(C : \mathcal{H} \to \mathcal{H}\) is linear with operator norm \(\|C\|_2 \leq L_C\).
  \item \(B_p(x) = U_p \Sigma_p V_p^\top x\) is linear with \(\|\Sigma_p\|_2 \leq L_{B_p}\).
  \item \(A_p : \mathcal{H} \to \mathcal{H}\) is 1-Lipschitz:
        \(\|A_p(u) - A_p(v)\| \leq \|u - v\|\) for all \(u,v\).
\end{enumerate}
Then the composite map \(\mathcal{O}_p := A_p \circ B_p \circ C\) is Lipschitz with
\[
  \|\mathcal{O}_p\|_{\mathrm{Lip}} \leq L_C \cdot L_{B_p}.
\]
\end{theorem}

\begin{proof}
For any \(x,y \in \mathcal{H}\),
\begin{align*}
  \|\mathcal{O}_p(x) - \mathcal{O}_p(y)\|
  &= \|A_p(B_p(Cx)) - A_p(B_p(Cy))\| \\
  &\leq \|B_p(Cx) - B_p(Cy)\|
     && \text{(1-Lipschitz \(A_p\))} \\
  &= \|B_p(C(x-y))\| \\
  &\leq \|B_p\|_2 \,\|C(x-y)\| \\
  &\leq \|B_p\|_2 \,\|C\|_2 \,\|x-y\|.
\end{align*}
Thus \(\|\mathcal{O}_p\|_{\mathrm{Lip}} \leq \|B_p\|_2 \,\|C\|_2 \leq L_C L_{B_p}\).
\end{proof}

\subsection*{A.2 Lipschitz Bounds for Resonance and Gates}

We now explicitly include the Lipschitz behavior of the resonance maps \(R_p(x)\) and the gates \(w_p(x)\).

Let \(x^{(\mathrm{gen})}\) and \(x^{(\mathrm{bio})}\) denote two modality-specific components of \(x\), obtained via fixed linear projections \(P_{\mathrm{gen}}, P_{\mathrm{bio}}\) with bounded norms:
\[
  x^{(\mathrm{gen})} = P_{\mathrm{gen}} x,
  \quad x^{(\mathrm{bio})} = P_{\mathrm{bio}} x,
  \quad
  \|P_{\mathrm{gen}}\|_2 \leq L_{\mathrm{gen}},
  \quad
  \|P_{\mathrm{bio}}\|_2 \leq L_{\mathrm{bio}}.
\]

Fix a prime \(p\) and its resonance matrix \(M_p\).
Define
\[
  R_p(x) =
  \frac{\langle x^{(\mathrm{gen})}, M_p x^{(\mathrm{bio})} \rangle}
       {\|x^{(\mathrm{gen})}\|\,\|M_p x^{(\mathrm{bio})}\| + \varepsilon},
\]
with \(\varepsilon > 0\), and
\[
  w_p(x) = \sigma(R_p(x) - \tau_p),
\]
where \(\sigma\) is the logistic sigmoid.

\begin{lemma}[Lipschitz Bound for Resonance and Gate]
Assume:
\begin{enumerate}[label=(\roman*),leftmargin=*]
  \item Inputs lie in a bounded set:
        \(\|x\| \leq R\) for all \(x\) in the domain of interest.
  \item \(M_p\) is bounded: \(\|M_p\|_2 \leq L_{M_p}\).
\end{enumerate}
Then there exist constants \(L_{R_p}, L_{w_p} > 0\) such that
\[
  |R_p(x) - R_p(y)| \leq L_{R_p} \,\|x-y\|,
  \quad
  |w_p(x) - w_p(y)| \leq L_{w_p} \,\|x-y\|
\]
for all \(x,y\) in the bounded set.
In particular, one may take
\[
  L_{R_p} =
    C_R \, L_{M_p}\,L_{\mathrm{gen}}\,L_{\mathrm{bio}},
  \quad
  L_{w_p} \leq \tfrac{1}{4} L_{R_p},
\]
for some constant \(C_R\) depending on \(R\) and \(\varepsilon\).
\end{lemma}

\begin{proof}[Proof Sketch]
The numerator of \(R_p(x)\) can be written as
\[
  N(x) = \langle P_{\mathrm{gen}} x, M_p P_{\mathrm{bio}} x \rangle.
\]
This is a smooth bilinear form in \(x\).
On a bounded set \(\|x\|\leq R\), one can bound the gradient norm
\(\|\nabla_x N(x)\|\) by a constant proportional to \(L_{M_p} L_{\mathrm{gen}} L_{\mathrm{bio}} R\).
Similarly, the denominator
\[
  D(x) = \|P_{\mathrm{gen}}x\|\,\|M_p P_{\mathrm{bio}}x\| + \varepsilon
\]
is bounded away from zero by \(\varepsilon\), and its gradient is bounded on the same set.
By the quotient rule and standard bounds for smooth functions on compact sets, we obtain a global Lipschitz constant \(L_{R_p}\) of the stated form.

For the gate, the logistic sigmoid \(\sigma\) satisfies \(|\sigma'(z)| \leq \tfrac{1}{4}\) for all \(z \in \mathbb{R}\).
Thus
\[
  |w_p(x) - w_p(y)| =
  |\sigma(R_p(x)-\tau_p) - \sigma(R_p(y)-\tau_p)|
  \leq \sup_z |\sigma'(z)|\,|R_p(x)-R_p(y)|
  \leq \tfrac{1}{4} L_{R_p} \|x-y\|.
\]
\end{proof}

\subsection*{A.3 Global Operator with Gate Lipschitzness}

We now refine the global Lipschitz bound for \(\mathcal{F}\) by incorporating the Lipschitz behavior of the gates \(w_p(x)\).

Recall:
\[
  \mathcal{F}(x) = \sum_{p \in P} w_p(x)\,\Pi_p\,\mathcal{O}_p(x).
\]

\begin{theorem}[Refined Global Lipschitz Bound]
Assume:
\begin{enumerate}[label=(\roman*),leftmargin=*]
  \item Conditions of the Channel Lipschitz Bound hold, with
        \(\|\mathcal{O}_p\|_{\mathrm{Lip}} \leq L_{\mathcal{O}_p}\).
  \item Each projector satisfies \(\|\Pi_p\|_2 \leq 1\).
  \item For each \(p\), the gate \(w_p\) is Lipschitz with constant \(L_{w_p}\), and
        \(|w_p(x)| \leq W_p\) on the domain of interest.
\end{enumerate}
Then for any \(x,y\),
\[
  \|\mathcal{F}(x) - \mathcal{F}(y)\|
  \leq
  \left(
    \sum_{p \in P} W_p L_{\mathcal{O}_p}
    + \sum_{p \in P} L_{w_p} C_{\mathcal{O}_p}
  \right)
  \|x-y\|,
\]
where \(C_{\mathcal{O}_p}\) is a uniform bound on \(\|\mathcal{O}_p(z)\|\) for \(z\) in the bounded set.
In particular, if the right-hand coefficient is less than 1, then \(\mathcal{F}\) is a contraction.
\end{theorem}

\begin{proof}
For any \(x,y\),
\begin{align*}
  \mathcal{F}(x) - \mathcal{F}(y)
  &= \sum_{p}
     \Big(
       w_p(x)\,\Pi_p\,\mathcal{O}_p(x)
       - w_p(y)\,\Pi_p\,\mathcal{O}_p(y)
     \Big) \\
  &= \sum_{p}
     \Big(
       w_p(x)\,\Pi_p(\mathcal{O}_p(x) - \mathcal{O}_p(y))
       + (w_p(x)-w_p(y))\,\Pi_p \mathcal{O}_p(y)
     \Big).
\end{align*}
Taking norms and using the triangle inequality,
\begin{align*}
  \|\mathcal{F}(x) - \mathcal{F}(y)\|
  &\leq \sum_p
    |w_p(x)|\,\|\Pi_p\|_2\,\|\mathcal{O}_p(x) - \mathcal{O}_p(y)\| \\
  &\quad
    + \sum_p
    |w_p(x) - w_p(y)|\,\|\Pi_p\|_2\,\|\mathcal{O}_p(y)\|.
\end{align*}
Using the bounds \(|w_p(x)| \leq W_p\), \(\|\Pi_p\|_2 \leq 1\), and \(\|\mathcal{O}_p\|_{\mathrm{Lip}} \leq L_{\mathcal{O}_p}\), we get
\[
  \|\mathcal{F}(x) - \mathcal{F}(y)\|
  \leq
  \sum_p W_p L_{\mathcal{O}_p} \|x-y\|
  + \sum_p L_{w_p} C_{\mathcal{O}_p} \|x-y\|,
\]
where we have assumed \(\|\mathcal{O}_p(z)\| \leq C_{\mathcal{O}_p}\) for all \(z\) in the bounded set.
This yields the stated bound.
\end{proof}

In practice, training is performed on bounded regions of state space and weights are regularized, making the constants finite and controllable.

\subsection*{A.4 Stability of the Full Recursion}

We now address the stability of the full recursion
\[
  x_{t+1} = x_t + \Lambda T(\mathcal{F}(x_t)).
\]

Define the map
\[
  G(x) := x + \Lambda T(\mathcal{F}(x)).
\]
We consider \(\Lambda T\) as a bounded linear operator on \(\mathcal{H}\).

\begin{theorem}[Contraction of the Full Recursion]
Assume:
\begin{enumerate}[label=(\roman*),leftmargin=*]
  \item \(\mathcal{F}\) is Lipschitz with constant \(L_{\mathcal{F}} < \infty\).
  \item \(T : \mathcal{H} \to \mathcal{H}\) is Lipschitz with constant \(L_T\).
  \item \(\Lambda : \mathcal{H} \to \mathcal{H}\) is linear with \(\|\Lambda\|_2 \leq \alpha\).
\end{enumerate}
Then for any \(x,y\),
\[
  \|G(x) - G(y)\|
  \leq (1 + \alpha L_T L_{\mathcal{F}})\,\|x-y\|.
\]
In particular:
\begin{itemize}[leftmargin=*]
  \item If \(\alpha L_T L_{\mathcal{F}} < 0\), the trivial bound is \(1\) (no contraction).
  \item If we modify the recursion to
        \(x_{t+1} = (1-\beta)x_t + \beta \Lambda T(\mathcal{F}(x_t))\)
        with \(\beta \in (0,1]\),
        then
        \[
          \|G_\beta(x) - G_\beta(y)\|
          \leq \bigl(1-\beta + \beta \alpha L_T L_{\mathcal{F}}\bigr)\,\|x-y\|,
        \]
        and \(G_\beta\) is a contraction whenever
        \(1-\beta + \beta \alpha L_T L_{\mathcal{F}} < 1\),
        i.e.
        \(\alpha L_T L_{\mathcal{F}} < 1\).
\end{itemize}
\end{theorem}

\begin{proof}
For the unmodified \(G(x) = x + \Lambda T(\mathcal{F}(x))\),
\begin{align*}
  \|G(x) - G(y)\|
  &= \|x-y + \Lambda T(\mathcal{F}(x)) - \Lambda T(\mathcal{F}(y))\| \\
  &\leq \|x-y\|
    + \|\Lambda\|_2 \,\|T(\mathcal{F}(x)) - T(\mathcal{F}(y))\| \\
  &\leq \|x-y\|
    + \alpha L_T \,\|\mathcal{F}(x) - \mathcal{F}(y)\| \\
  &\leq \big(1 + \alpha L_T L_{\mathcal{F}}\big)\,\|x-y\|.
\end{align*}
Thus the Lipschitz constant is at most \(1 + \alpha L_T L_{\mathcal{F}}\), which is \(\geq 1\) unless the second term is negative (not the case here).

For the modified recursion
\[
  G_\beta(x) = (1-\beta)x + \beta \Lambda T(\mathcal{F}(x)),
\]
we get
\begin{align*}
  \|G_\beta(x) - G_\beta(y)\|
  &\leq (1-\beta)\|x-y\|
    + \beta \alpha L_T L_{\mathcal{F}} \|x-y\| \\
  &= \bigl(1-\beta + \beta \alpha L_T L_{\mathcal{F}}\bigr)\,\|x-y\|.
\end{align*}
Hence \(G_\beta\) is a contraction if and only if
\(1-\beta + \beta \alpha L_T L_{\mathcal{F}} < 1\),
equivalently \(\alpha L_T L_{\mathcal{F}} < 1\).
\end{proof}

\begin{corollary}[Unique Stable Fixed Point for the Relaxed Recursion]
Under the assumptions of the previous theorem, if
\[
  \alpha L_T L_{\mathcal{F}} < 1,
\]
then for any \(\beta \in (0,1]\), the relaxed recursion
\[
  x_{t+1} = (1-\beta)x_t + \beta \Lambda T(\mathcal{F}(x_t))
\]
has a unique fixed point \(x^\star\), and for any initialization \(x_0\), the sequence \(x_t\) converges to \(x^\star\).
\end{corollary}

\begin{proof}
The map \(G_\beta\) is a contraction on the complete metric space \(\mathcal{H}\) with Lipschitz constant strictly less than 1, so the Banach fixed-point theorem applies.
\end{proof}

\subsection*{A.5 Interpretation}

In practice, \(\Lambda\) and \(T\) are under the designer's control (e.g.\ \(\Lambda\) as a step-size / gain and \(T\) as a simple linear or low-Lipschitz head), while \(\mathcal{F}\) is trained to satisfy a Lipschitz bound \(L_{\mathcal{F}}\) through spectral normalization and regularization.
The condition
\[
  \alpha L_T L_{\mathcal{F}} < 1
\]
thus becomes a \emph{design objective}: choose architecture and training such that the effective loop gain remains below 1, guaranteeing stability and a unique equilibrium for the recursive dynamics.

\nocite{*}
\bibliographystyle{plainnat}
\bibliography{references}
\end{document}